\documentclass[10pt,a4paper]{amsart}
\usepackage[margin=19mm]{geometry}
\usepackage{amsmath,amssymb,mathtools}
\usepackage[scr=rsfs,cal=euler]{mathalfa}
\usepackage{xfrac}
\usepackage[unicode,hidelinks,bookmarks=false]{hyperref}
\newcommand{\ie}{i.e.\ }
\newcommand{\cf}{cf.\ }
\newcommand{\resp}{respectively\ }
\newcommand{\Subsection}{\subsection}
\providecommand{\nobreakdash}{\nobreak\hbox{-}}
\setlength{\emergencystretch}{2em}
\multlinegap0pt
\usepackage{amssymb}
\usepackage{mathtools}
\newtheorem{theorem}{Theorem}
\newcommand{\ID}{{\mathbb D}}
\newcommand{\abs}[1]{\left\vert#1\right\vert}
\title{Integral mean estimates for $(\alpha,\beta)$-harmonic functions}
\author{Zhi-Gang Wang, Brindha Valson E, R. Vijayakumar}
\date{Source: arXiv:2603.11449v1 (CC BY 4.0)}
\begin{document}
\maketitle

\noindent\emph{Source and licence.} Integral mean estimates for $(\alpha,\beta)$-harmonic functions. Version arXiv:2603.11449v1 (CC BY 4.0), distributed by arXiv under the Creative Commons Attribution 4.0 International licence, http://creativecommons.org/licenses/by/4.0/, as recorded in the arXiv licence metadata for this exact version and shown on the abstract page https://arxiv.org/abs/2603.11449v1. Full licence notice: \texttt{licenses/arxiv-2603.11449-CC-BY-4.0.txt}.\par\smallskip

\noindent\textit{Editorial scope.} Counted unit: the article's theorem t4 (source0.tex lines 673--680). The article's complete original proof of it is delivered with it (source0.tex lines 681--755, one \texttt{proof} environment of 75 lines). No mathematics is written, repaired or re-derived: the statement and the proof are the article's own lines, in the article's own notation, and no retained line is substituted or re-flowed.  Retained uncounted as the source-local context the proof needs: lines 233--235; lines 551--559.  Nothing was omitted from any retained span, so the package carries no omission record.  No retained line contains a citation, so the wrapper carries no bibliography and no bibliography entry.  No retained line contains a figure environment, an image-inclusion command or drawing code, so no image is delivered and the package declares no asset dependency.  Editorial formatting only, in addition: an AMS \texttt{amsart} preamble that restates the article's own declarations and macros used by the retained lines, namely the environments \texttt{theorem} on separate counters, together with the standard packages those lines need.  The retained environments print in this document as Theorem 1 in order of appearance, whereas the article prints the same items with the numbering of its own full document.  The complete native source of the article as distributed in the arXiv e-print for this exact version is bundled unchanged as \texttt{sources/196216/source0.tex}.
		\begin{equation}\label{poisson_integral}
			w(z) = K_{\alpha,\beta}[f](z) = \frac{1}{2\pi} \int_0^{2\pi} K_{\alpha,\beta}\left(z e^{-it}\right) f\left(e^{it}\right)dt,
		\end{equation}

		\begin{equation}\label{series of (alpha,beta)}
			\begin{aligned}
				u(z)
				&= \sum_{m=0}^{\infty}
				c_m\,F(-\alpha,m-\beta; m+1; |z|^2)\,z^m  \\
				&\quad\ \ +\sum_{m=1}^{\infty}
				c_{-m}\,F(-\beta,m-\alpha; m+1; |z|^2)\,\bar z^{\,m}
			\end{aligned}
			\end{equation}

\begin{theorem}\label{t4}
	Let $w(z)=K_{\alpha,\beta}[f](z)$ be an $(\alpha,\beta)$-harmonic function with $\mbox{\rm Re}(\alpha+\beta)>-1$ defined on  $\mathbb{D}$  with $f\in L^p(\mathbb{T}),1\leq p\leq \infty$. If $w(z)$ is of the form \eqref{series of (alpha,beta)}, then 
	\begin{itemize}
		\item [(i)] $|c_k|\leq |c_{\alpha,\beta}|\dfrac{(\alpha+1)_{k}}{k!}\|f\|_{L^p(\mathbb{T})},\ \ |c_{-k}|\leq|c_{\alpha,\beta}|\dfrac{(\beta+1)_{k}}{k!}\|f\|_{L^p(\mathbb{T})}.$
		\item [(ii)]$\dfrac{k!}{(\beta+1)_{k}}|c_k|+\dfrac{k!}{(\alpha+1)_{k}}|c_{-k}|\leq 2|c_{\alpha,\beta}|C_q\|f\|_{L^p(\mathbb{T})}$.
	\end{itemize}
In particular, when $p=\infty$, $C_q=\frac{2}{\pi}$ and $\alpha=\beta=0$, we have $$|c_k|+|c_{-k}|\leq\frac{4}{\pi}\|f\|_{\infty}.$$
\end{theorem}

\begin{proof}
	Let $w(z)$ be an $(\alpha,\beta)$-harmonic function defined on $\ID$. By \eqref{poisson_integral},
    we have
	\begin{equation}\label{47}
			w(z) = \frac{1}{2\pi} \int_0^{2\pi} c_{\alpha,\beta} \frac{\left(1 - |z|^2\right)^{\alpha + \beta + 1}}{(1 - z e^{-it})^{\alpha+1} (1 - \overline{z} e^{it})^{\beta+1}} f\left(e^{it}\right)dt.
	\end{equation}
	Now, by direct calculation, we get
	
	\begin{align*}
		w_{z}(0)
		&= \frac{c_{\alpha,\beta}}{2\pi}\int_{0}^{2\pi} (\alpha+1)\,e^{it} f(e^{it})\,dt,\\
		w_{ z}^{(2)}(0)
		&= \frac{c_{\alpha,\beta}}{2\pi}\int_{0}^{2\pi} (\alpha+2)(\alpha+1)\,e^{2it} f(e^{it})\,dt,\\ 
		&\ \vdots \notag\\
		w_{ z}^{(k)}(0)
		&= \frac{c_{\alpha,\beta}}{2\pi}\int_{0}^{2\pi} (\alpha+1)_{k}\, e^{ikt} f(e^{it})\,dt.
		\end{align*}
	Similarly, we have
	$$w_{ \overline{z}}^{(k)}(0)
	= \frac{c_{\alpha,\beta}}{2\pi}\int_{0}^{2\pi} (\beta+1)_{k}\, e^{ikt} f(e^{it})\,dt.
	$$
Thus, for \(k=1,2,\ldots\), we know that 
	$$
	c_k
	= \frac{w_z^{(k)}(0)}{k!}
	= \,
	\frac{c_{\alpha,\beta}}{2\pi}\int_{0}^{2\pi} \frac{(\alpha+1)_{k}}{k!}\,e^{ikt} f\left(e^{it}\right)dt.
	$$
	Moreover, we find that
	$$
	c_{-k}
	= \frac{w_{\bar z}^{(k)}(0)}{k!}
	= \,
	\frac{1}{2\pi}\int_{0}^{2\pi} \frac{(\beta+1)_{k}}{k!}\,e^{ikt} f\left(e^{it}\right)\,dt.
	$$
	Applying H\"older's inequality, we get
	$$|c_k|\leq |c_{\alpha,\beta}|\frac{(\alpha+1)_{k}}{k!}\|f\|_{L^p(\mathbb{T})}\ \ \text{and}\ \ |c_{-k}|\leq|c_{\alpha,\beta}|\frac{(\beta+1)_{k}}{k!}\|f\|_{L^p(\mathbb{T})}.$$
Let $c_k=|c_k|e^{i\alpha_k}$, $c_{-k}=|c_{-k}|e^{-i\beta_k}$ and $\theta_k=\dfrac{\alpha_k+\beta_k}{2}$. Then 
$$
|c_k|=\frac{(\alpha+1)_{k}}{ k!}\frac{c_{\alpha,\beta}}{2\pi}
\int_{0}^{2\pi}
f\left(e^{it}\right)e^{-ikt}e^{-i\alpha_k}dt$$ and
$$
|c_{-k}|=\frac{(\beta+1)_{k}}{ k!}
\frac{c_{\alpha,\beta}}{2\pi}\int_{0}^{2\pi}
f\left(e^{it}\right)e^{ikt}e^{i\beta_k}dt.
$$
Therefore, we deduce that
\begin{align*}
\begin{split}
	&\frac{ k!}{(\alpha+1)_{k}}\,|c_k|+ \frac{k!}{(\beta+1)_{k}}\,|c_{-k}|
	\\=&\left|
	\frac{c_{\alpha,\beta}}{2\pi i}\int_{0}^{2\pi}f\left(e^{it}\right) \left(e^{-ikt}e^{-i\alpha_k}+e^{ikt}e^{i\beta_k}\right)dt\right|\\ \leq&\frac{|c_{\alpha,\beta}|}{2\pi}
	\int_{0}^{2\pi}\abs{f\left(e^{it}\right)}\,\bigl|e^{-ikt}e^{-i\alpha_k}+e^{ikt}e^{i\beta_k}\bigr|\,dt
	\\=&\frac{|c_{\alpha,\beta}|}{\pi}\int_{0}^{2\pi}\abs{f\left(e^{it}\right)}\,\bigl|\cos k(t+\theta_k)\bigr|dt\\ \leq&
	2|c_{\alpha,\beta}|\left(\frac{1}{2\pi}\int_{0}^{2\pi}
	\abs{f\left(e^{it}\right)}^{p}\,dt\right)^{\frac{1}{p}}
	\left(\frac{1}{2\pi}\int_{0}^{2\pi}
	|\cos k(t+\theta_k)|^{q}\,dt\right)^{\frac{1}{q}}\\=&
	2|c_{\alpha,\beta}|C_q\,\|f\|_{L^p(\mathbb{T})} ,
    \end{split}
	\end{align*}
where
\[
C_q
=
\left(
\frac{1}{2\pi}
\int_{0}^{2\pi}
|\cos(kt)|^{q}\,dt
\right)^{\frac{1}{q}},
\]
and \(q\) is the conjugate exponent of \(p\), satisfying
\(\frac{1}{p}+\frac{1}{q}=1\).
\end{proof}

\end{document}
